% ---------------------------------------------------------------
% preamble/theorems.tex - theorem-like environments in coloured boxes
% Same names as usual: \begin{definition}[title] \label{def:x} ... \end{definition}
% All numbered ones share one counter per section: Theorem 1.1, Definition 1.2, ...
% Starred forms (theorem*, ...) are unnumbered.
% ---------------------------------------------------------------

\numberwithin{equation}{section}

% No extra space inside the boxes; the boxes set the spacing
\declaretheoremstyle[
  spaceabove=0pt, spacebelow=0pt, headfont=\bfseries\sffamily,
  bodyfont=\itshape, headpunct={}, postheadspace=0.6em,
]{boxplain}
\declaretheoremstyle[
  spaceabove=0pt, spacebelow=0pt, headfont=\bfseries\sffamily,
  bodyfont=\normalfont, headpunct={}, postheadspace=0.6em,
]{boxdef}

\declaretheorem[style=boxplain, name=\txttheorem, numberwithin=section]{theorem}
\declaretheorem[style=boxplain, name=\txtlemma,       sibling=theorem]{lemma}
\declaretheorem[style=boxplain, name=\txtproposition, sibling=theorem]{proposition}
\declaretheorem[style=boxplain, name=\txtcorollary,   sibling=theorem]{corollary}
\declaretheorem[style=boxplain, name=\txtconjecture,  sibling=theorem]{conjecture}
\declaretheorem[style=boxplain, name=\txttheorem, numbered=no]{theorem*}
\declaretheorem[style=boxplain, name=\txtlemma,   numbered=no]{lemma*}

\declaretheorem[style=boxdef, name=\txtdefinition, sibling=theorem]{definition}
\declaretheorem[style=boxdef, name=\txtexample,    sibling=theorem]{example}
\declaretheorem[style=boxdef, name=\txtexercise,   sibling=theorem]{exercise}
\declaretheorem[style=boxdef, name=\txtdefinition, numbered=no]{definition*}
\declaretheorem[style=boxdef, name=\txtexample,    numbered=no]{example*}

\declaretheorem[style=boxdef, name=\txtremark, sibling=theorem]{remark}
\declaretheorem[style=boxdef, name=\txtremark, numbered=no]{remark*}

% \begin{solution} ... \end{solution}  proof layout with "Solution" title
\newenvironment{solution}
  {\begin{proof}[\txtsolution]}
  {\end{proof}}

% --- Box colours: blue definitions, red theorems, green examples, grey remarks ---
\tcbset{
  thmbox/.style={
    enhanced, breakable, sharp corners, frame hidden,
    left=7pt, right=7pt, top=5pt, bottom=5pt,
    before skip=0.9em, after skip=0.9em,
  },
}
\tcolorboxenvironment{definition}{thmbox, colback=blue!4, borderline west={3pt}{0pt}{blue!60!black}}
\tcolorboxenvironment{definition*}{thmbox, colback=blue!4, borderline west={3pt}{0pt}{blue!60!black}}
\tcolorboxenvironment{theorem}{thmbox, colback=red!4, borderline west={3pt}{0pt}{red!60!black}}
\tcolorboxenvironment{theorem*}{thmbox, colback=red!4, borderline west={3pt}{0pt}{red!60!black}}
\tcolorboxenvironment{lemma}{thmbox, colback=red!3, borderline west={3pt}{0pt}{red!45!black}}
\tcolorboxenvironment{lemma*}{thmbox, colback=red!3, borderline west={3pt}{0pt}{red!45!black}}
\tcolorboxenvironment{proposition}{thmbox, colback=red!3, borderline west={3pt}{0pt}{red!45!black}}
\tcolorboxenvironment{corollary}{thmbox, colback=red!3, borderline west={3pt}{0pt}{red!45!black}}
\tcolorboxenvironment{conjecture}{thmbox, colback=orange!5, borderline west={3pt}{0pt}{orange!70!black}}
\tcolorboxenvironment{example}{thmbox, colback=green!4, borderline west={3pt}{0pt}{green!45!black}}
\tcolorboxenvironment{example*}{thmbox, colback=green!4, borderline west={3pt}{0pt}{green!45!black}}
\tcolorboxenvironment{exercise}{thmbox, colback=black!3, borderline west={3pt}{0pt}{black!50}}
\tcolorboxenvironment{remark}{thmbox, colback=black!3, borderline west={3pt}{0pt}{black!35}}
\tcolorboxenvironment{remark*}{thmbox, colback=black!3, borderline west={3pt}{0pt}{black!35}}
